% Rubric: Background and Research Aims + Originality and Contribution.
% Establish the problem and its significance, state a focused research
% question or modelling objective, position the work against the literature,
% and clearly separate the team's contribution from existing or taught models.

\section{Introduction}
\label{sec:introduction}

\subsection{Background and research aims}
A station closure can sever a railway corridor before any further facility
fails. Substitute buses may reconnect journeys, but the paths they restore can
also concentrate traffic at surviving rail facilities. We investigate these
competing effects in the Transperth network: infrastructure connectivity,
secondary overload failures and the fraction of terminal pairs that remain
served. Fewer open rail facilities need not imply less service when another
transport layer is available.

\citet{albert2000error} established the distinction between random errors and
targeted removal in heterogeneous networks. Their results motivate the
random-versus-targeted comparison here; they do not establish that Transperth
shares their degree distribution or resilience. Our audited railway
has 86 mapped stations and 85 adjacent segments and is a tree, so it lacks
alternative rail paths around a removed station. Its inputs represent
Monday 5 October 2026, 07:00--09:00 in Australia/Perth, with 271 selected rail
trips and 4,324 selected bus trips. These are scheduled services, not measured
passenger journeys. Showgrounds remains a mapped station although it has no
weekday stop event in this window.

\citet{motter2002cascade} showed how a local attack can cause further overload
failures when shortest-path loads are recomputed against capacities fixed from
the intact network. We use that capacity construction and compare it with local
redistribution, in which a failed facility forwards its carried load to
survivors; the comparison tests whether an apparent tolerance transition
depends on the update mechanism. In the bus layer, restored paths
can increase surviving rail loads even while improving terminal reachability.
Apparent avalanche scaling requires a separate statistical test: following
\citet{clauset2009power}, a visual slope alone is not evidence for a power law.

Our research questions are: \textbf{RQ1}, how do random and targeted removal
fragment the railway, and which station characteristics identify bottlenecks?
\textbf{RQ2}, how do capacity tolerance and load-update rules affect the size
and duration of secondary failures? \textbf{RQ3}, when do bus backup links
improve served travel, and can their rerouting aggravate a rail cascade?
\textbf{RQ4}, how sensitive are these conclusions to sampling, tolerance-grid
resolution and the assumed load/capacity model, and is avalanche scaling
statistically supported?

\subsection{Contribution}
Our contribution is a reproducible transport case study that joins an audited
GTFS-derived topology to removal, overload and recovery experiments. We expand
express-service traversals over mapped adjacent segments rather than treating
skipped stops as shortcuts. We separate persistent passenger terminals from
closable rail facilities, retain evidence for substitute links, and compare
infrastructure damage with service outcomes. Paired scenarios and bootstrap
intervals expose uncertainty without treating scheduled departures as observed
demand. The capacity law and robustness questions come from existing work;
the audited Transperth construction, explicit comparison of load-update
assumptions, and evidence-linked recovery analysis are the team's extension.
The implementation is network simulation; the unit's taught agent-based models
are not reused.
